\documentclass[12pt]{article}
\usepackage{mathpazo}
\usepackage[T1]{fontenc}
\usepackage[utf8]{inputenc}
\usepackage{geometry}
\geometry{verbose,tmargin=1cm,bmargin=2.7cm,lmargin=2.7cm,rmargin=2cm}
\pagestyle{empty}
\usepackage{setspace}
\usepackage{amsmath,amssymb}
\usepackage{booktabs}
\usepackage{tabularx}
\usepackage{enumitem}
\setlist[enumerate]{leftmargin=1cm}
\setlist[itemize]{leftmargin=0.5cm}
\onehalfspacing
\setlength{\parskip}{\smallskipamount}
\setlength{\parindent}{0pt}
\setlength{\emergencystretch}{3em}
\date{}
\author{}
\newcommand{\solution}[1]{}
\newcommand{\hint}[1]{\par\smallskip{\small\textit{Hint:} #1}\par}

\begin{document}

\title{\textsc{Assignment I: Trickle-Down vs Capital Gains}}
\maketitle
\vspace{-3mm}\thispagestyle{empty}

\textbf{Objective:} This project teaches you to solve a heterogeneous agent model with a thick upper tail of wealth and imperfect competition, and to use it to study how rising labor income inequality shapes the distribution of wealth under perfect and imperfect competition, in steady state and along transitions.
\begin{itemize}
\item \textbf{Problem:} The problem consists of
\begin{enumerate}
\item A number of questions (page 2)
\item A model (page 3 onward, incl. calibration and numerical implementation)
\item Some code files you can start from
\end{enumerate}
\item \textbf{Code:} The problem is designed to be solved with the \texttt{GEModelTools} package.
\item \textbf{Figures and tables:} Please use the provided functions to plot figures and tables.
\item \textbf{Structure:} Your project should consist of
\begin{enumerate}
\item A single self-contained pdf-file with all results
\item A single Jupyter notebook showing how the results are produced
\item Well-documented \texttt{.py} files
\end{enumerate}
\item \textbf{Hand-in:} Upload a single zip-file on Absalon (and nothing else)
\item \textbf{Deadline:} 26 October 2026
\item \textbf{Exam:} This project will be a part of your exam portfolio.
You can incorporate feedback before handing in the final version.
\end{itemize}

\newpage

\section*{Questions}

\begin{enumerate}
\item \textbf{Calibration.}
\begin{enumerate}
\item \emph{Pareto tails.} For wealthy households, the saving rate out of income tends to one when $\Sigma<\sigma$, and wealth and consumption have Pareto tails with coefficients (the derivation is in Section 4)
$$\zeta_a=-\frac{\log(1-\xi)}{\log(1+r)},\qquad \zeta_c=\frac{\sigma}{\Sigma}\zeta_a.$$
Explain the intuition behind each formula. Gaillard et al.\ (2024) estimate $\zeta_a=1.35$ and $\zeta_c=3.02$ in US data. With $r=3\%$ and $\sigma=2$, find $\xi$ and $\Sigma$.
\solution{1a}
\item \emph{Household problem.} Complete the inversion of the Euler equation (\texttt{c\_endo} and \texttt{m\_endo}) in \texttt{solve\_hh\_backwards} in \texttt{household\_problem.py}.
\hint{The Euler equation of a surviving household is in Section 4. Death is an exogenous state, so the expectation step \texttt{z\_trans@v\_a} already includes the survival probability: the marginal value of $a_{-}$ is zero for a newborn, since its $a_{-}$ is the wealth of the parent who died, which goes to the bequest pool.}
\solution{1b}
\item \emph{Wealth distribution.} Complete \texttt{blocks\_ss} and \texttt{find\_ss\_calibrate} in \texttt{steady\_state.py}. Calibrate $\beta$, $\gamma$ and $\bar a$ such that the asset market clears and the model matches the bottom 50\% and top 10\% wealth shares in 1970 (\texttt{wealth\_psz.csv}).
\hint{Given $Y=1$, $r$ and $W/Y$, the formula for Tobin's Q in question 2(b) gives $K$. Then use \texttt{obj\_ss}, which returns the asset market clearing error for a given $K$, inside a root-finder over $(\beta,\log\gamma,\log\bar a)$, starting from $\beta=0.9$, $\gamma=2$ and $\bar a=5$; you should find $\beta\approx0.84$, $\gamma\approx1.92$ and $\bar a\approx5.48$.}
\solution{1c}
\item \emph{Policy functions.} For wealthy households ($a_{-}\gg\bar a$ and labor income negligible), consumption and net savings are approximately
$$c\approx Ba_{-}^{\Sigma/\sigma},\qquad B=\left[\frac{\gamma(1+r)^{-\Sigma}}{1-\beta(1-\xi)(1+r)^{1-\Sigma}}\right]^{-1/\sigma},\qquad a-a_{-}\approx ra_{-}$$
(the derivation is in Section 4).
Plot the consumption function and net savings $a-a_{-}$ of the calibrated model as functions of $a_{-}$, in logs, together with these approximations.
Why do the rich save so much in this model? How does this depend on $\Sigma$ relative to $\sigma$?
How would the policy functions look for high values of wealth in a homothetic model without taste for wealth ($\gamma=0$)?
\solution{1d}
\item \emph{Perfect competition.} Construct an economy with $\mu=1$ that keeps all the other parameters of the $\mu=1.2$ economy, including $\Gamma$, $\delta$, $\gamma$ and $\bar a$, except $\beta$, which you re-solve such that $r=3\%$ (complete \texttt{find\_ss\_beta} in \texttt{steady\_state.py}). How do $Y$, $w$, $K/Y$, $W/Y$, $Q$, $\beta$ and the 1970 wealth distribution differ from the economy with $\mu=1.2$? Why?
\hint{$K$ follows from the investment condition at $r=3\%$. Then find $\beta$ with \texttt{optimize.brentq} on $[0.5,0.95]$, using \texttt{obj\_ss} at this $K$.}
\solution{1e}
\end{enumerate}

\item \textbf{Steady-state impact of rising permanent income inequality.}
\begin{enumerate}
\item Define the stationary equilibrium of the model in Section 1. Show that the goods market clears when the asset market clears (Walras' law).
\solution{2a}
\item \emph{Tobin's Q.} Show that in steady state
$$Q\equiv\frac{q}{K}=(1-\tau_d)\left[1+\alpha\left(1-\frac{1}{\mu}\right)\frac{Y/K}{r}\right].$$
Interpret the two terms. What is $Q$ under perfect competition, $\mu=1$?
\solution{2b}
\item \emph{Linear approximation.} Let $\partial\log A$ be the partial equilibrium change in log asset demand when the permanent income grid moves from 1970 to 2019 (\texttt{labor\_psz.csv}) and households face the 1970 prices, $\varepsilon^A=\partial\log A/\partial r$ the semi-elasticity of asset demand (with $w$ and bequests moving with $r$ as in steady state), $\varepsilon^q=-\partial\log(q+B)/\partial r$ the semi-elasticity of asset supply and $\varepsilon^K=-\partial\log K/\partial r$.
\begin{enumerate}[label=(\roman*)]
\item Show that, to first order,
$$dr=-\frac{\partial\log A}{\varepsilon^A+\varepsilon^q},\qquad d\log w=-\alpha\varepsilon^K dr,\qquad d\log\frac{W}{Y}=\frac{\varepsilon^q-\alpha\varepsilon^K}{\varepsilon^A+\varepsilon^q}\partial\log A.$$
\hint{Write the steady-state asset market clearing condition as a function of $r$ and the permanent income grid, $\log A(r,\{z_i\})=\log\left(q(r)+B\right)$, where $w$ and the bequest are themselves functions of $r$. Then take the total derivative.}
\item Interpret the formulas.
\item Compute $\partial\log A$ and the elasticities in the model for $\mu=1.2$ and $\mu=1$, and use the formulas to predict the changes in $r$, $w$ and $W/Y$. Why do markups dampen the fall in $r$?
\end{enumerate}
\solution{2c}
\item \emph{Rising permanent income inequality.}
\begin{enumerate}[label=(\roman*)]
\item Replace the 1970 permanent income grid by the 2019 grid and compute the new steady state of both economies, $\mu=1.2$ and $\mu=1$.
\item Compare the changes in $r$, $w$, $K/Y$, $W/Y$, $Q=q/K$ and the wealth shares in the two economies. How close is the linear approximation of question 2(c)?
\item Solve the household problem with the 1970 permanent income grid and the 2019 prices (markets do not have to clear). Compare the average wealth of the top 1\% and of the top 0.1\% relative to output (the 2019 output, which goes with the 2019 prices), and the top 1\% wealth share, to 1970. What is the impact of the change in prices on wealth inequality under perfect and under imperfect competition? Why?
\end{enumerate}
\hint{For (iii), copy the 2019 model, give it the 1970 grid, and call \texttt{solve\_hh\_ss} and \texttt{simulate\_hh\_ss}.}
\solution{2d}
\end{enumerate}

\item \textbf{One-time shock.} Assume that the rise in permanent income inequality is an unexpected permanent shock in 1970, after which households and firms have perfect foresight.
\begin{enumerate}
\item \emph{Firm value.} Complete \texttt{production\_firm} and the computation of $q_t$ in \texttt{mutual\_fund} in \texttt{blocks.py}. Why is $q_t$ computed backwards in time?
\hint{Use the no-arbitrage condition in Section 1. Start from the terminal condition $q_T=q_{ss}$, the firm value in the new steady state, and iterate backwards to $t=0$. \texttt{next(q,t,ss.q)} returns $q_{t+1}$, and $q_{ss}$ in the last period.}
\solution{3a}
\item Compute the transition from the 1970 to the 2019 steady state for $\mu=1.2$ and $\mu=1$. Plot the return on wealth, average Tobin's Q ($q_t/K_t$), the wealth-output ratio and the top 1\% wealth share, in deviations from 1970.
Explain the return on wealth in the first period. Why is there no capital gain under perfect competition?
\solution{3b}
\item \emph{(Hard)} Along the transition, decompose the change since 1970 in the average wealth of the top 1\% and of the top 0.1\%, relative to output, into
\begin{itemize}
\item a partial equilibrium effect: households have the 2019 permanent income grid and face the 1970 prices (divide by 1970 output),
\item a general equilibrium effect: households have the 1970 permanent income grid and face the realized paths of $r^B_t$, $cg_t$, $w_t$ and $b_t$ (divide by realized output).
\end{itemize}
Compare the general equilibrium effect under $\mu=1$ and $\mu=1.2$. Does its sign change over time?
\hint{Use \texttt{decompose\_hh\_path}, starting from the 1970 distribution (\texttt{Dbeg}). For the partial equilibrium effect, call it on a copy of the 1970 model with the 2019 grid and \texttt{use\_inputs=[]}; for the general equilibrium effect, call it on the 1970 model with \texttt{use\_inputs} set to the four prices and their realized paths in \texttt{custom\_paths}.}
\solution{3c}
\item \emph{Conclusion.} Given your answers to questions 2 and 3, discuss who benefits from the change in prices caused by the rise in permanent income inequality, under perfect and under imperfect competition, in the short run and in the long run.
\solution{3d}
\end{enumerate}

\item \textbf{Myopic expectations (optional).}
Assume instead that the permanent income grid moves every 5 years along a linear path from the 1970 to the 2019 grid, in 10 steps, and that each time households and firms believe the new grid is permanent.
Compute the resulting path as a sequence of unexpected shocks, and compare the return on wealth, average Tobin's Q and the average wealth of the top 0.1\% to question 3.
\hint{Each step is a transition from the current state to the steady state households now believe in. Five years after a shock, the next one starts from \texttt{ini = \{'Dbeg':path.Dbeg[5], 'K':path.K[4,0], 'q':path.q[4,0], 'r':path.r[5,0]\}}: the bonds held in year 5 were bought in year 4 and pay the return $r_5$ expected then.}
\solution{4}
\end{enumerate}


\newpage

\section{Model}

\paragraph*{Households.}
The model is annual. There is a continuum of households with permanent labor productivity $z_i$ from five groups (bottom 50\%, next 40\%, top 10--1\%, top 1--0.1\% and top 0.1\% of labor income), and idiosyncratic productivity $e$ with $\log e'=\rho_e\log e+\epsilon'$.
Households supply one unit of labor inelastically, with efficiency $z_i e$, and aggregate efficiency units are normalized to one.
Each period a household dies with probability $\xi$ and is replaced by a newborn of the same permanent type, whose idiosyncratic productivity continues from its parent's.
Households value consumption and end-of-period wealth:
\begin{align*}
v_t(z_i,e,a_{-}) & =\max_{c,a}\frac{c^{1-\sigma}}{1-\sigma}+\gamma\frac{(a+\bar a)^{1-\Sigma}}{1-\Sigma}+\beta(1-\xi)\mathbb{E}\left[v_{t+1}(z_i,e',a)\right]\\
\text{s.t. } & a+c=(1+r^B_t)a_{-}+cg_t\,s(a_{-})a_{-}+(1-\tau_l)w_t z_i e,\qquad a\geq0.
\end{align*}
Since $\Sigma<\sigma$, wealth is a luxury good, and $\bar a$ makes the taste for wealth weaker for poor households.
A newborn has no wealth of its own and receives the bequest $b_t z_i$ instead of $(1+r^B_t)a_{-}+cg_t\,s(a_{-})a_{-}$.

\paragraph*{Bequests.}
The wealth of the dead is taxed at rate $\tau_E$ and the rest is given to newborns in proportion to their permanent productivity. Since death is independent of wealth, the dead hold a share $\xi$ of wealth and $b_t=(1-\tau_E)(1+r^a_t)A_{t-1}$.

\paragraph*{Final-good producer.}
A competitive final-good producer combines labor $L_t$ and a CES bundle of intermediate goods $x_{j,t}$, $j\in[0,1]$:
$$Y_t=\Gamma L_t^{1-\alpha}X_t^{\alpha},\qquad X_t=\left(\int_0^1 x_{j,t}^{1/\mu}\,dj\right)^{\mu},$$
where $\mu\geq1$, so that the elasticity of substitution between varieties is $\mu/(\mu-1)$. The final good is the numéraire.
Maximizing $Y_t-w_tL_t-\int_0^1 p_{j,t}x_{j,t}\,dj$ gives the wage $w_t=(1-\alpha)Y_t/L_t$ and the demand for each variety
$$x_{j,t}=X_t\left(\frac{p_{j,t}}{P^X_t}\right)^{-\frac{\mu}{\mu-1}},\qquad P^X_t=\alpha\frac{Y_t}{X_t},$$
so that spending on intermediate goods is $\int_0^1 p_{j,t}x_{j,t}\,dj=\alpha Y_t$. The final-good producer makes no profits.

\paragraph*{Intermediate producers.}
Each intermediate producer $j$ owns its capital stock $k_{j,t-1}$, chosen in the previous period, and produces $x_{j,t}=k_{j,t-1}$. It sets its price $p_{j,t}$ given the demand curve above, invests, and pays the dividend
$$d_{j,t}=p_{j,t}x_{j,t}-\left[k_{j,t}-(1-\delta)k_{j,t-1}\right].$$
Its ex-dividend value $q_{j,t}$ satisfies the no-arbitrage condition
$$q_{j,t-1}(1+r_t)=(1-\tau_d)d_{j,t}+q_{j,t},$$
where $r_t$ is the after-tax return on wealth between $t-1$ and $t$, and shareholders pay the dividend tax $\tau_d$. The producer chooses its price and its investment to maximize its value, taking the demand curve and $P^X_t$ as given. You do not need to derive the first-order conditions, which are:
\begin{itemize}
\item \emph{Pricing.} The price is a markup $\mu$ over the marginal cost of capital (the user cost),
$$p_{j,t}=\mu\left(r_t+\delta\right).$$
\item \emph{Investment.} One unit of the final good invested in period $t$ yields the marginal revenue $p_{j,t+1}/\mu$ and the undepreciated capital $1-\delta$ in period $t+1$, so
$$1+r_{t+1}=\frac{p_{j,t+1}}{\mu}+1-\delta.$$
\end{itemize}
The dividend tax does not enter the investment condition, since it scales all dividends by the same factor.

\paragraph*{Symmetric equilibrium.}
All intermediate producers make the same choices, so $x_{j,t}=X_t=K_{t-1}$ and $p_{j,t}=P^X_t=\alpha Y_t/K_{t-1}$. With $L_t=1$,
$$Y_t=\Gamma K_{t-1}^{\alpha},\qquad w_t=(1-\alpha)Y_t,\qquad r_t=\frac{\alpha}{\mu}\frac{Y_t}{K_{t-1}}-\delta.$$
Intermediate producers earn $\alpha Y_t$, of which the rents are $\alpha(1-1/\mu)Y_t$, and aggregate dividends and firm value $q_t=\int_0^1 q_{j,t}\,dj$ satisfy
$$d_t=\alpha Y_t-I_t,\qquad I_t=K_t-(1-\delta)K_{t-1},\qquad q_{t-1}(1+r_t)=(1-\tau_d)d_t+q_t.$$
Under perfect competition ($\mu=1$), intermediate producers earn no rents and $q_t=(1-\tau_d)K_t$.

\paragraph*{Mutual fund.}
A mutual fund holds all equity and government bonds $B$, $A_t=q_t+B$. Absent surprises, all wealth earns $r^B_t=r_t$.
When the economy is hit by an unexpected shock in period 0, bonds pay the return expected when they were bought, $r^B_0=r_{-1}$, while equity earns an unexpected capital gain.
This gain is paid to households in proportion to their equity exposure $s(a_{-})a_{-}$, where the equity share of wealth follows the portfolio rule estimated on the SCF by Elina and Huleux,
$$s(a)=\min\left\{\frac{\psi_0}{1+\exp\left(-\psi_1\left(\log(a/A_{1970})-\psi_2\right)\right)},1\right\},$$
where $A_{1970}$ is aggregate wealth in the 1970 steady state.
The capital gain per unit of exposure is
$$cg_0=\frac{(1-\tau_d)d_0+q_0-(1+r_{-1})q_{-1}}{\int s(a_{-})a_{-}\,d\boldsymbol{D}_0},$$
and $cg_t=0$ for $t\geq1$. Households' return on wealth is $r^B_t+cg_t\,s(a_{-})$, so the rich, who hold more equity, receive a larger share of the gain.
The average return on wealth is $r^a_t$, with $1+r^a_0=\left[(1-\tau_d)d_0+q_0+(1+r_{-1})B\right]/(q_{-1}+B)$.

\paragraph*{Government.}
Debt is constant and spending adjusts:
$$G_t=\tau_l w_t+\tau_d d_t+\tau_E\,\xi(1+r^a_t)A_{t-1}-r^B_t B.$$

\paragraph*{Market clearing.}
\begin{enumerate}
\item Asset market: $A_t=q_t+B=A^{hh}_t$
\item Goods market: $Y_t=C^{hh}_t+I_t+G_t$
\end{enumerate}

\section{Calibration}

\begin{table}[h]
\centering
\caption{Parameters}
\small
\begin{tabularx}{\textwidth}{lXll}
\toprule
 & Parameter & Value & Source \\
\midrule
$\sigma$ & CRRA on consumption & 2.0 & \\
$\xi$, $\Sigma$ & death probability, curvature of taste for wealth & & question 1(a) \\
$\beta$, $\gamma$, $\bar a$ & discount factor, taste for wealth & & question 1(c) \\
$z_i$ & permanent productivity & & PSZ labor income shares \\
$\rho_e$, $\sigma_e$ & persistence and std. of $\log e$ & 0.95, 0.545 & Storesletten et al. \\
$\alpha$ & capital share & 1/3 & \\
$\mu$ & markup & 1.2 & \\
$\Gamma$, $\delta$ & technology, depreciation & & $Y=1$, $r=3\%$, $W/Y=3.7$ \\
$B$ & government debt & 0.19 & \\
$\tau_l$ & labor income tax & 0.30 & \\
$\tau_E$ & estate tax & 0.30 & \\
$\tau_d$ & dividend tax & 0.30 & \\
$\psi_0$, $\psi_1$, $\psi_2$ & portfolio rule & 1.0, 0.374, 2.598 & Elina and Huleux, SCF \\
\bottomrule
\end{tabularx}
\end{table}

To calibrate the 1970 steady state, set $Y=1$, $r=3\%$ and $W/Y=3.7$.
This gives $q=W-B$, $K$ from Tobin's Q, $\Gamma$ from output and $\delta$ from the investment condition.
Then solve for $\beta$, $\gamma$ and $\bar a$ (question 1(c)).
The economy with $\mu=1$ keeps all these parameters except $\beta$ (question 1(e)).

\section{Numerical implementation}

\begin{itemize}
\item \textbf{States.} The five permanent types are fixed states (\texttt{Nfix = 5}). Idiosyncratic productivity is discretized with the Rouwenhorst method with 4 states.
Death is handled as an extra exogenous state: the household is either alive or newborn, $4\times2=8$ states. A newborn does not use $a_{-}$ in its budget.
\item \textbf{Asset grid.} 500 points between 0 and $10^{10}$, finest around $a=5$. Because the grid spans many orders of magnitude, rounding errors in savings at the top of the grid are about $10^{-7}$, so the tolerance when solving the household problem is $10^{-6}$.
\item \textbf{Distribution.} The top of the wealth distribution converges slowly, so the tolerance when simulating the distribution is very tight ($10^{-17}$).
\item \textbf{Transitions.} The only unknown is $K$, with asset market clearing as the target. Use $T=300$. A transition from an old steady state needs \texttt{ini} to contain \texttt{Dbeg}, \texttt{K}, \texttt{q} and \texttt{r}. GEModelTools prints a warning that some terminal values differ slightly from the steady state; this is expected, since the top of the distribution converges slowly.
\item \textbf{Helper functions.} \texttt{tails.py} computes wealth shares and the average wealth of top groups.
\end{itemize}

\section{Derivations (optional)}

This section derives the Euler equation and the approximations used in questions 1(a), 1(b) and 1(d). You are not asked to reproduce them.

\paragraph*{Policy functions.} In steady state ($r^B=r$, $cg=0$) the Euler equation of a surviving household is
$$c_t^{-\sigma}=\beta(1-\xi)(1+r)c_{t+1}^{-\sigma}+\gamma(a_t+\bar a)^{-\Sigma},$$
where $1-\xi$ appears because wealth is worth nothing to a household that dies. For a wealthy household, $a_t\gg\bar a$ and labor income is negligible. Guess that $c=Ba_{-}^{\Sigma/\sigma}$. If $\Sigma<\sigma$, consumption relative to wealth, $c/a_{-}=Ba_{-}^{\Sigma/\sigma-1}$, goes to zero, so the budget constraint gives $a_t\approx(1+r)a_{-}$ and net savings are $a_t-a_{-}\approx ra_{-}$. Next period's consumption is then $c_{t+1}=Ba_t^{\Sigma/\sigma}=B(1+r)^{\Sigma/\sigma}a_{-}^{\Sigma/\sigma}$. Substituting into the Euler equation,
$$B^{-\sigma}a_{-}^{-\Sigma}=\beta(1-\xi)(1+r)B^{-\sigma}(1+r)^{-\Sigma}a_{-}^{-\Sigma}+\gamma(1+r)^{-\Sigma}a_{-}^{-\Sigma}.$$
The powers of $a_{-}$ cancel, which verifies the guess, and
$$B=\left[\frac{\gamma(1+r)^{-\Sigma}}{1-\beta(1-\xi)(1+r)^{1-\Sigma}}\right]^{-1/\sigma},$$
which requires $\beta(1-\xi)(1+r)^{1-\Sigma}<1$.

\paragraph*{Pareto tail of wealth.} At the top, the wealth of a surviving household grows at the gross rate $1+r$, and each period a household dies with probability $\xi$ and its dynasty restarts from the bequest. A household that has survived $n$ periods since the restart has wealth $x\propto(1+r)^n$, and the share of households that have survived at least $n$ periods is $(1-\xi)^n$. With $n=\log x/\log(1+r)$,
$$P(a>x)\propto(1-\xi)^{\log x/\log(1+r)}=x^{\log(1-\xi)/\log(1+r)},$$
so wealth has a Pareto tail with coefficient $\zeta_a=-\log(1-\xi)/\log(1+r)$.

\paragraph*{Pareto tail of consumption.} Since $c=Ba^{\Sigma/\sigma}$ at the top,
$$P(c>y)=P\left(a>(y/B)^{\sigma/\Sigma}\right)\propto y^{-\zeta_a\sigma/\Sigma},$$
so consumption has a Pareto tail with coefficient $\zeta_c=\frac{\sigma}{\Sigma}\zeta_a$.

\section*{References}
\begin{itemize}
\item Elina, E. and R. Huleux. \emph{From Income to Wealth Inequality: Trickle-Down vs Capital Gains}. Working paper.
\item Gaillard, A., C. Hellwig, P. Wangner and N. Werquin (2024). \emph{Consumption, Wealth, and Income Inequality: A Tale of Tails}. Working paper.
\item Piketty, T., E. Saez and G. Zucman (2018). Distributional National Accounts: Methods and Estimates for the United States. \emph{Quarterly Journal of Economics} 133(2), 553--609.
\item Storesletten, K., C. Telmer and A. Yaron (2004). Consumption and Risk Sharing over the Life Cycle. \emph{Journal of Monetary Economics} 51(3), 609--633.
\end{itemize}

\end{document}
